\documentclass[10pt,twocolumn]{article}

% ============================================================
%  Packages
% ============================================================
\usepackage[utf8]{inputenc}
\usepackage[T1]{fontenc}
\usepackage{mathptmx}
\usepackage{amsmath,amssymb,amsfonts}
\usepackage{graphicx}
\usepackage{booktabs}
\usepackage{multirow}
\usepackage{hyperref}
\usepackage[margin=0.75in,top=1in,bottom=1in]{geometry}
\usepackage[numbers,sort&compress]{natbib}
\usepackage{algorithm}
\usepackage{algorithmic}
\usepackage{xcolor}
\usepackage{caption}
\usepackage{subcaption}
\usepackage{float}
\usepackage{microtype}
\usepackage{enumitem}
\usepackage{setspace}
\usepackage{authblk}
\usepackage{xurl}

\graphicspath{{../figures/}}

\setcounter{topnumber}{3}
\setcounter{bottomnumber}{2}
\setcounter{totalnumber}{5}
\renewcommand{\topfraction}{0.9}
\renewcommand{\bottomfraction}{0.8}
\renewcommand{\textfraction}{0.1}
\renewcommand{\floatpagefraction}{0.8}
\setcounter{dbltopnumber}{2}
\renewcommand{\dbltopfraction}{0.9}

\hypersetup{
    colorlinks=true,
    linkcolor=blue!70!black,
    citecolor=blue!70!black,
    urlcolor=blue!70!black
}

\captionsetup{font=small,labelfont=bf}
\setlist{nosep,leftmargin=*}

\usepackage{titlesec}
\titleformat{\section}{\large\bfseries}{\thesection}{1em}{}
\titleformat{\subsection}{\normalsize\bfseries}{\thesubsection}{0.8em}{}
\titleformat{\subsubsection}{\normalsize\itshape}{\thesubsubsection}{0.8em}{}
\titlespacing*{\section}{0pt}{1.5ex plus 0.5ex minus 0.2ex}{0.8ex plus 0.2ex}
\titlespacing*{\subsection}{0pt}{1.2ex plus 0.3ex minus 0.2ex}{0.5ex plus 0.1ex}

% ============================================================
%  Title
% ============================================================
\title{\Large\bfseries Modular Versus Monolithic Architectures for Gene Regulatory
Network Edge Prediction: Gradient Stability Analysis and a Controlled
Cross-Architecture Comparison}

\author{Evint Leovonzko}

\date{}

% ============================================================
\begin{document}
\maketitle
\thispagestyle{empty}

% ============================================================
%  Abstract
% ============================================================
\begin{abstract}
\noindent\textbf{Motivation:}
Neural network architectures for gene regulatory network (GRN) inference differ
fundamentally in how they represent the interaction between a transcription
factor (TF) and a target gene: \emph{modular} two-tower models score pairs via
a learned similarity function in a shared latent space, while \emph{monolithic}
cross-encoders concatenate all pairwise features before a shared MLP.
Whether architectural modularity benefits or hinders performance on this
pairwise classification task remains under-explored, particularly when
confounded by training instabilities that can mask genuine model differences.

\noindent\textbf{Results:}
We present a controlled comparison of a two-tower hybrid embedding model and a
cross-encoder on a curated human brain single-cell RNA-seq dataset of 47{,}388
TF--gene candidate pairs.
A prerequisite diagnostic investigation identifies three compounding gradient
failures in the original two-tower training implementation---a loss-function
double-sigmoid, gradient replacement (rather than accumulation) in FC layer
backward passes, and catastrophic gradient clipping under temperature
scaling---each of which, individually and in combination, collapses training to
near-chance accuracy.
After correcting these failures with per-batch Adam and a numerically stable
gradient estimator, the two-tower model achieves 80.90\% accuracy
(95\%~CI:~[80.62\%, 82.36\%], AUROC~0.810, 5 seeds) under balanced sampling.
The cross-encoder reaches 83.03\% accuracy (95\%~CI:~[81.78\%, 83.57\%],
AUROC~0.904) under the same conditions, a statistically meaningful gap of
+2.1~pp accuracy and +0.094~AUROC.
Under realistic imbalanced sampling (5:1 negative ratio), the gap widens
further: cross-encoder degrades modestly ($-1.6$~pp) while the two-tower model
degrades substantially ($-6.9$~pp), indicating that the modular architecture's
cosine similarity scoring is sensitive to class-prior shift at the decision
boundary.
All models are implemented in pure Rust without external deep learning
dependencies.

\noindent\textbf{Availability:}
Source code, result JSONs, and analysis scripts are available at
\url{https://github.com/Evintkoo/module-regularized-grn}.

\noindent\textbf{Keywords:} gene regulatory networks, two-tower architecture,
cross-encoder, modular networks, gradient analysis, single-cell RNA-seq,
negative sampling
\end{abstract}

\vspace{0.5em}

% ============================================================
\section{Introduction}
\label{sec:introduction}

Gene regulatory networks (GRNs) encode the transcription factor (TF)--target
gene interactions that control cell fate and tissue homeostasis~\cite{davidson2002genomic,barabasi2004network}.
Predicting which TFs regulate which genes from single-cell RNA-seq data is a
binary classification problem over the combinatorial space of TF--gene
pairs~\cite{pratapa2020benchmarking,chen2019evaluating}.

A fundamental design choice in neural GRN classifiers is how to represent the
\emph{interaction} between a TF and a gene.
Two broad families exist:

\textbf{Modular (two-tower) models} encode TFs and genes through separate
parametric pathways into a shared latent space, then score interactions via a
lightweight similarity function (dot product, cosine).
This design is well-studied in retrieval and recommendation
systems~\cite{yi2019sampling} but its application to GRN inference has received
limited rigorous evaluation.

\textbf{Monolithic (cross-encoder) models} concatenate all pairwise features
before processing them through a unified MLP.
Cross-encoders have no inductive bias toward symmetric or geometrically
interpretable representations, but provide greater capacity for modeling
arbitrary nonlinear feature interactions.

Comparing these families requires identical data, evaluation protocol, and
implementation quality.
In practice, published comparisons are often confounded by optimizer differences,
unequal hyperparameter tuning, or training instabilities that are not reported.

In this work, we present a fully controlled comparison of a two-tower and a
cross-encoder on the same dataset with the same optimizer.
A prerequisite contribution is a careful gradient analysis that identifies three
compounding implementation failures in the original two-tower training script
(\S\ref{sec:grad_analysis}), together with a per-batch Adam training procedure
that corrects them.
This analysis is of independent value: the identified failure modes (loss
double-sigmoid, backward-pass gradient replacement, and temperature-amplified
clipping) can occur undetected in any from-scratch neural network implementation.

Our contributions are:
\begin{enumerate}
    \item We identify and diagnose three gradient failures in a from-scratch
    two-tower implementation that each collapse training to near-chance accuracy,
    and provide a corrected training procedure (\S\ref{sec:training}).
    \item We propose a \textbf{cross-encoder architecture} (Model~D) that
    concatenates TF and gene embeddings with their element-wise product and
    cell-type expression features, processing the joint representation through
    a three-layer MLP to a raw logit (\S\ref{sec:cross_encoder}).
    \item We conduct a \textbf{controlled cross-architecture comparison} over
    five seeds $\times$ two negative-sampling ratios, demonstrating that the
    monolithic cross-encoder consistently outperforms the two-tower model on
    AUROC (+0.094 at 1:1; +0.172 at 5:1) and is substantially more robust to
    negative-ratio shift (\S\ref{sec:results}).
\end{enumerate}

\subsection{Related work}
\label{sec:related}

\subsubsection{Classical GRN inference}
GENIE3~\cite{huynh2010inferring} and GRNBoost2~\cite{moerman2019grnboost2}
use tree-based regression.
SCENIC~\cite{aibar2017scenic} augments regression with motif enrichment.
ARACNE~\cite{margolin2006aracne} and Bayesian network methods~\cite{friedman2000bayesian}
apply mutual information and probabilistic graphical models.
These methods typically achieve AUROC 0.65--0.75 on standardized
benchmarks~\cite{pratapa2020benchmarking}.

\subsubsection{Deep learning for GRN inference}
GNNLink~\cite{wang2023gnnlink} treats GRN inference as graph link prediction.
Transformer-based methods scGREAT~\cite{chen2024scgreat} and
TRENDY~\cite{yuan2025trendy} apply self-attention over expression profiles.
Foundation models scGPT~\cite{cui2024scgpt} and
Geneformer~\cite{theodoris2023transfer} pretrain on millions of cells.
These GPU-requiring methods report accuracies of 85--90\%.

\subsubsection{Knowledge graph embeddings and two-tower retrieval}
Our two-tower model adapts ideas from TransE~\cite{bordes2013translating} and
two-tower retrieval~\cite{yi2019sampling} to the GRN pairwise scoring task.
Cross-encoder models in our setting are analogous to the reranker stage in
dual-encoder/reranker retrieval pipelines, where a cross-encoder refines
candidate scores by processing query--document pairs jointly.

% ============================================================
\section{Materials and Methods}
\label{sec:methods}

\subsection{Problem formulation}
\label{sec:problem}

Let $\mathcal{T} = \{t_1, \ldots, t_M\}$ be a set of transcription factors
and $\mathcal{G} = \{g_1, \ldots, g_N\}$ a set of target genes.
We learn $s_\theta(t, g) \in [0, 1]$ estimating the probability of a regulatory
edge $t \to g$, formulated as supervised binary classification over candidate pairs.

\subsection{Data sources and preprocessing}
\label{sec:data}

\subsubsection{Single-cell expression data}
Human brain single-cell RNA-seq data (H5AD format) provides cell-type expression
profiles $\mathbf{x}_g \in \mathbb{R}^{C}$, $C = 11$ brain cell types, where
$x_{g,c}$ is the mean log-normalized expression of gene $g$ in cell type $c$.
In the current evaluation, 0\% of TFs and genes in the prior databases map to
a non-zero expression vector (``0\% expression coverage''), so the model
effectively operates on 512-dimensional embeddings alone; the 11-dimensional
expression channels are zero-padded but preserved in the architecture for
forward compatibility.

\subsubsection{Prior regulatory knowledge}
Positive regulatory pairs are compiled from DoRothEA~\cite{garcia2019benchmark}
and TRRUST v2~\cite{han2018trrust}, yielding 29{,}084 unique TF--gene pairs
spanning 1{,}061 TFs and 10{,}247 genes.
After intersection with the scRNA-seq gene space, 23{,}694 positive pairs
survive, yielding a dataset of 47{,}388 total examples at 1:1 positive-to-negative ratio.

\subsubsection{Negative sampling}
\label{sec:negative}
Negative examples are sampled uniformly from $\mathcal{T} \times \mathcal{G}$
excluding the positive set.
We evaluate two negative-sampling regimes:
\begin{itemize}
    \item \textbf{1:1 balanced}: equal positives and negatives in the training
    set (standard benchmark setting).
    \item \textbf{5:1 imbalanced}: five negatives per positive in the training
    set, approximating realistic GRN sparsity.
\end{itemize}
Validation and test sets are always 1:1 balanced for comparability across
experiments.

\subsubsection{Data partitioning}
The dataset is partitioned 70/15/15 (train/val/test) by stratified random
sampling.
No TF--gene pair appears across splits; individual TFs and genes may appear
paired with different partners.

\subsection{Model A: Two-Tower (HybridEmbeddingModel)}
\label{sec:architecture}

\subsubsection{Overview}
Each TF $t_i$ and gene $g_j$ is associated with a learnable embedding:
$\mathbf{e}^{\text{TF}}_i, \mathbf{e}^{\text{Gene}}_j \in \mathbb{R}^{d_e}$,
$d_e = 512$.
The hybrid input for each entity concatenates its embedding with its expression
profile:
\begin{equation}
    \mathbf{h}_0^{(\ell)} = \bigl[\mathbf{e}^{(\ell)} \mathbin{;} \mathbf{x}^{(\ell)}\bigr]
    \in \mathbb{R}^{d_e + C} = \mathbb{R}^{523}
\end{equation}

\subsubsection{Encoder towers}
Two-layer feedforward encoders (one per entity type, separate parameters):
\begin{align}
    \mathbf{h}_1^{(\ell)} &= \text{ReLU}\!\bigl(\mathbf{W}_1^{(\ell)}\mathbf{h}_0^{(\ell)} + \mathbf{b}_1^{(\ell)}\bigr), \quad
    \mathbf{W}_1^{(\ell)} \in \mathbb{R}^{d_h \times 523} \\
    \mathbf{z}^{(\ell)} &= \mathbf{W}_2^{(\ell)}\mathbf{h}_1^{(\ell)} + \mathbf{b}_2^{(\ell)}, \quad
    \mathbf{W}_2^{(\ell)} \in \mathbb{R}^{d_o \times d_h}
\end{align}
with $d_h = 512$, $d_o = 512$.

\subsubsection{Interaction scoring}
\begin{equation}
    s_\theta(t_i, g_j) = \sigma\!\left(\frac{\langle \mathbf{z}^{\text{TF}}_i,\,\mathbf{z}^{\text{Gene}}_j \rangle}{\tau}\right),
    \qquad \tau = 0.05
\end{equation}
The sigmoid output is the final prediction (probability); total parameters $\approx 5.2\mathrm{M}$.

\subsection{Model D: Cross-Encoder}
\label{sec:cross_encoder}

\subsubsection{Overview}
Unlike the two-tower design, the cross-encoder processes TF and gene information
\emph{jointly} through a single MLP.
Given TF embedding $\mathbf{e}^{\text{TF}}_i \in \mathbb{R}^{d_e}$, gene
embedding $\mathbf{e}^{\text{Gene}}_j \in \mathbb{R}^{d_e}$, and their
respective expression profiles, the input is:
\begin{equation}
    \mathbf{u}_{ij} = \bigl[\mathbf{e}^{\text{TF}}_i \mathbin{;} \mathbf{e}^{\text{Gene}}_j \mathbin{;}
    \mathbf{e}^{\text{TF}}_i \odot \mathbf{e}^{\text{Gene}}_j \mathbin{;}
    \mathbf{x}^{\text{TF}}_i \mathbin{;} \mathbf{x}^{\text{Gene}}_j\bigr]
    \label{eq:cross_input}
\end{equation}
where $\odot$ denotes element-wise product.
Input dimensionality: $3 d_e + 2C = 3 \times 512 + 2 \times 11 = 1{,}558$.

The element-wise product term $\mathbf{e}^{\text{TF}} \odot \mathbf{e}^{\text{Gene}}$
encodes coordinate-wise multiplicative interactions between TF and gene
embedding dimensions, providing a low-rank feature-product expansion without
requiring a quadratic parameter count.
This term is analogous to the multiplicative gate in factored bilinear models~\cite{yang2015embedding}.

\subsubsection{Architecture}
Three fully connected layers process the joint representation to a raw logit:
\begin{align}
    \mathbf{h}_1 &= \text{ReLU}(\mathbf{W}_1 \mathbf{u}_{ij} + \mathbf{b}_1), \quad
    \mathbf{W}_1 \in \mathbb{R}^{d_h \times 1558} \\
    \mathbf{h}_2 &= \text{ReLU}(\mathbf{W}_2 \mathbf{h}_1 + \mathbf{b}_2), \quad
    \mathbf{W}_2 \in \mathbb{R}^{d_h \times d_h} \\
    \hat{y} &= \mathbf{w}_3^\top \mathbf{h}_2 + b_3 \in \mathbb{R}
    \label{eq:cross_logit}
\end{align}
with $d_h = 512$.
The output is a raw \emph{logit}; BCE loss is applied externally via
$\sigma(\hat{y})$.
Total parameters $\approx 5.6\mathrm{M}$.

\subsubsection{Key structural differences from two-tower}
Table~\ref{tab:arch_comparison} summarizes the architectural differences.
The cross-encoder sacrifices the parameter efficiency and latency advantages
of separate towers in favor of richer pairwise feature representation.
In particular, the $\mathbf{e}^{\text{TF}} \odot \mathbf{e}^{\text{Gene}}$ term
gives the model direct access to multiplicative interactions that the two-tower
model must reconstruct implicitly through dot-product scoring.

\begin{table}[!htbp]
\centering
\caption{Structural comparison of the two architectures. TT: Two-Tower; CE: Cross-Encoder.}
\label{tab:arch_comparison}
\small
\begin{tabular}{@{}lcc@{}}
\toprule
\textbf{Property} & \textbf{TT} & \textbf{CE} \\
\midrule
Entity paths        & Separate & Joint \\
Interaction features & Dot product & $\mathbf{e}_i \odot \mathbf{e}_j$ (explicit) \\
Output              & Sigmoid score & Raw logit \\
Score function      & $\sigma(\langle\mathbf{z}_i,\mathbf{z}_j\rangle/\tau)$ & $\sigma(\hat{y})$ \\
Input dim. (FC)     & 523 per tower & 1{,}558 joint \\
Total FC params     & ${\sim}$0.8M & ${\sim}$1.1M \\
Total params        & ${\sim}$5.2M & ${\sim}$5.6M \\
Symmetric by design & Yes & No \\
\bottomrule
\end{tabular}
\end{table}

\subsection{Gradient stability analysis and optimizer selection}
\label{sec:grad_analysis}

During development, the two-tower model trained with an epoch-level gradient
update procedure converged to near-chance accuracy (${\sim}$50--59\%) across
all seeds despite correct forward pass logic.
We identified three compounding failure modes through systematic gradient flow
analysis.

\subsubsection{Failure 1: Double-sigmoid in the loss backward pass}
\label{sec:failure1}
The existing \texttt{bce\_loss\_backward} utility function applies sigmoid
internally before computing $(p - y)/n$:
\begin{equation}
    \texttt{bce\_loss\_backward}(\mathbf{x}) = \frac{\sigma(\mathbf{x}) - \mathbf{y}}{n}
\end{equation}
When the model's forward pass already outputs $s = \sigma(\cdot/\tau)$, passing
$s$ to this function computes $\sigma(s)$ a second time---a double-sigmoid.
For any prediction $s \in (0, 1)$, $\sigma(s) \in (0.5, 0.73)$, compressing
all gradient signals into a narrow regime and reversing the gradient direction
near the tails.

\subsubsection{Failure 2: Gradient replacement in batch iteration}
\label{sec:failure2}
The \texttt{LinearLayer::backward} method in the Rust implementation replaces
(not accumulates) stored gradients on each call:
\begin{equation}
    \nabla \mathbf{W} \leftarrow \mathbf{X}^\top \cdot \delta \quad
    \text{(replaces on every batch)}
\end{equation}
With an epoch-level gradient update (accumulate all batches, then step), this
means only the \emph{last} mini-batch's gradient survives into the parameter
update.
For $N_{\text{train}} = 33{,}171$ examples and batch size 256, the final batch
contains ${\sim}147$ examples.
The effective FC-layer learning signal per epoch is $\approx 0.44\%$ of the
full dataset.
Embeddings, updated via index-based gradient accumulation ($+=$), do not
exhibit this deficiency and continue to learn correctly.

\subsubsection{Failure 3: Gradient annihilation from temperature-amplified clipping}
\label{sec:failure3}
The temperature parameter $\tau = 0.05$ amplifies gradients flowing through
the scoring function by a factor of $1/\tau = 20$:
\begin{equation}
    \frac{\partial \mathcal{L}}{\partial \mathbf{z}} \propto \frac{p - y}{\tau \cdot n}
\end{equation}
With the stable gradient estimator described in \S\ref{sec:stable_grad}, typical
gradient norms reach ${\sim}1{,}766$ per batch.
A gradient clipping threshold of $\delta_{\max} = 1.0$ therefore scales all
gradients by $1.0 / 1766 \approx 0.0006$---effectively zeroing all parameter
updates.

\subsubsection{Corrected training procedure}
\label{sec:training}

We resolve all three failures simultaneously:

\paragraph{Stable gradient estimator.}
\label{sec:stable_grad}
Rather than passing model outputs through \texttt{bce\_loss\_backward},
we compute a numerically stable per-sample gradient directly:
\begin{equation}
    \delta_i = \frac{p_i - y_i}{n}, \qquad p_i = s_\theta(t_i, g_i) \in (0, 1)
    \label{eq:stable_grad}
\end{equation}
This avoids $1/(p(1{-}p))$ blow-up for confident predictions and the
double-sigmoid distortion.
Internally, \texttt{model.backward} multiplies by $\partial\sigma/\partial z =
p(1{-}p)/\tau$, yielding a net per-sample gradient of
$(p_i{-}y_i)\cdot p_i(1{-}p_i)/(n\tau)$---numerically bounded and correct in direction.

\paragraph{Per-batch Adam with increased clip threshold.}
Each mini-batch triggers a full Adam parameter update:
$(\mathbf{W}^{(\ell)}_k, \mathbf{b}^{(\ell)}_k, \mathbf{e}^{\text{TF}}, \mathbf{e}^{\text{Gene}})$
are updated after every batch, not every epoch.
This ensures all FC-layer parameters receive gradients from the full training
signal.
We use $\eta = 10^{-3}$, $\beta_1 = 0.9$, $\beta_2 = 0.999$, $\epsilon = 10^{-8}$,
and gradient clip threshold $\delta_{\max} = 5.0$, consistent with a $\sim$353$\times$
increase in effective clip budget relative to the original $\delta_{\max} = 1.0$.

\paragraph{Weight initialization.}
Embedding and weight matrices are initialized from $\mathcal{N}(0, \sigma^2)$
with $\sigma = 0.01$, reducing initial prediction variance and stabilizing
early-epoch gradients.

The corrected implementation is used for all two-tower experiments in this paper.
The cross-encoder uses the same per-batch Adam procedure with the same hyperparameters,
applied directly to the logit output (no temperature scaling required).

\paragraph{Impact of each failure mode.}
Table~\ref{tab:failure_modes} quantifies the isolated and compound effects of
each failure on held-out accuracy.

\begin{table}[!htbp]
\centering
\caption{Accuracy under different combinations of gradient failures (1:1 ratio,
seed 42, reported on test set after 60 epochs). Epoch-level: epoch-level SGD;
double-$\sigma$: bce\_loss\_backward called on sigmoid output; clip=1.0: gradient
clip threshold 1.0.}
\label{tab:failure_modes}
\small
\begin{tabular}{@{}lccc@{}}
\toprule
\textbf{Configuration} & \textbf{Acc.} & \textbf{AUROC} & \textbf{F1} \\
\midrule
Epoch-level SGD + double-$\sigma$ + clip=1.0 & 79.53\% & --- & --- \\
Epoch-level SGD (LR=0.005) & 58.82\% & --- & --- \\
Per-batch Adam + double-$\sigma$ + clip=1.0 & 49.81\% & 0.5045 & 0.000 \\
\textbf{Per-batch Adam + stable grad + clip=5} & \textbf{80.93\%} & \textbf{0.8093} & \textbf{0.8156} \\
\bottomrule
\end{tabular}
\end{table}

Epoch-level SGD with double-sigmoid achieves 79.53\% Val.\ Acc.\ at epoch 10
but the model was never evaluated on the test set under those conditions, and
train/val loss diverges (train $\approx 0.87$, val $\approx 3.10$) indicating
severe overfitting.
Per-batch Adam with the broken double-sigmoid and tight clip collapses to
chance (49.81\%), illustrating that clip=1.0 annihilates all gradient signal
under $\tau = 0.05$.
The corrected procedure achieves 80.93\% on the test set in a numerically
stable and interpretable training run.

\subsection{Ensemble construction}
\label{sec:ensemble_method}

For both models, we train $K = 5$ independent instances with seeds
$\{42, 123, 456, 789, 1337\}$, differing only in initialization and batch
ordering.
Ensemble predictions are formed by averaging individual sigmoid outputs:
\begin{equation}
    s_{\text{ens}}(t, g) = \frac{1}{K}\sum_{k=1}^{K} s_{\theta_k}(t, g)
\end{equation}
For the cross-encoder, individual logit outputs are passed through sigmoid
before averaging.

\subsection{Evaluation protocol}
\label{sec:evaluation}

We report accuracy, F1, and AUROC on the 1:1 balanced test set ($n = 7{,}109$).
AUROC is computed via trapezoidal integration over the sorted score vector.
F1 uses the default threshold of 0.5.
Bootstrap 95\% confidence intervals ($B = 1{,}000$ resamples, seed 42) are
computed for accuracy.
We report mean $\pm$ standard deviation across 5 seeds as a measure of training
stability.

\subsection{Implementation}
\label{sec:implementation}

The complete pipeline is implemented in Rust (edition 2021) using
\texttt{ndarray} v0.15, \texttt{ndarray-rand}, and \texttt{hdf5} for H5AD
parsing.
No external deep learning frameworks are used.
Adam optimizer states (first and second moment estimates) are managed manually,
one Adam state struct per model.
The cross-encoder's per-batch Adam is applied to embeddings and all three FC
layers; the two-tower's per-batch Adam is applied to both embedding tables and
both towers' two FC layers each.

% ============================================================
\section{Results}
\label{sec:results}

\subsection{Main comparison: two-tower versus cross-encoder}
\label{sec:main_results}

Table~\ref{tab:main_comparison} presents the primary comparison across all
architectures and negative-sampling ratios.
Both models are evaluated on the same 1:1 balanced test set.

\begin{table}[!htbp]
\centering
\caption{Primary comparison. Acc: mean test accuracy $\pm$ SD across 5 seeds.
CI: 95\% bootstrap CI. AUROC, F1: mean across 5 seeds. Ens: 5-model ensemble
accuracy. All test sets are 1:1 balanced ($n = 7{,}109$).}
\label{tab:main_comparison}
\small
\begin{tabular}{@{}llccccc@{}}
\toprule
\textbf{Model} & \textbf{Ratio} & \textbf{Acc ($\pm$SD)} & \textbf{CI [95\%]} & \textbf{AUROC} & \textbf{F1} & \textbf{Ens} \\
\midrule
Two-Tower    & 1:1 & 80.90 $\pm$ 0.59 & [80.6, 82.4] & 0.810 & 0.807 & 83.47 \\
Two-Tower    & 5:1 & 74.03 $\pm$ 1.54 & [74.4, 76.3] & 0.743 & 0.660 & 75.92 \\
\midrule
Cross-Encoder & 1:1 & 83.03 $\pm$ 0.48 & [81.8, 83.6] & 0.904 & 0.831 & 84.02 \\
Cross-Encoder & 5:1 & 81.47 $\pm$ 0.74 & [79.8, 81.6] & 0.915 & 0.783 & 82.33 \\
\bottomrule
\end{tabular}
\end{table}

\textbf{At balanced (1:1) sampling}, the cross-encoder outperforms the two-tower
by +2.13~pp in accuracy (83.03\% vs.\ 80.90\%) and by +0.094 in AUROC
(0.904 vs.\ 0.810).
Both improvements are consistent across all five seeds (Table~\ref{tab:per_seed}).
The cross-encoder's variance is also lower ($\sigma = 0.48$ vs.\ $0.59$~pp),
indicating a more stable optimization landscape.
The F1 gap (+0.024) is smaller than the accuracy gap, reflecting that the
two-tower achieves comparable recall but somewhat lower precision.

\textbf{At imbalanced (5:1) sampling}, the gap widens dramatically in AUROC
(0.915 vs.\ 0.743, $\Delta = +0.172$) and substantially in accuracy
(81.47\% vs.\ 74.03\%, $\Delta = +7.44$~pp).
The two-tower's accuracy drops $-6.87$~pp when moving from 1:1 to 5:1 training,
versus only $-1.56$~pp for the cross-encoder.
This differential sensitivity to negative-ratio shift is analyzed in \S\ref{sec:ratio_analysis}.

\begin{table}[!htbp]
\centering
\caption{Per-seed results under balanced (1:1) sampling.}
\label{tab:per_seed}
\small
\begin{tabular}{@{}cccccc@{}}
\toprule
\multirow{2}{*}{\textbf{Seed}} & \multicolumn{2}{c}{\textbf{Two-Tower}} & & \multicolumn{2}{c}{\textbf{Cross-Encoder}} \\
\cmidrule{2-3} \cmidrule{5-6}
 & Acc.\ (\%) & AUROC & & Acc.\ (\%) & AUROC \\
\midrule
42   & 80.93 & 0.809 & & 83.63 & 0.906 \\
123  & 79.80 & 0.799 & & 83.04 & 0.905 \\
456  & 81.50 & 0.816 & & 82.66 & 0.903 \\
789  & 80.94 & 0.811 & & 83.47 & 0.908 \\
1337 & 81.33 & 0.813 & & 82.35 & 0.898 \\
\midrule
Mean $\pm$ SD & 80.90 $\pm$ 0.59 & 0.810 & & 83.03 $\pm$ 0.48 & 0.904 \\
\bottomrule
\end{tabular}
\end{table}

\subsection{Effect of negative sampling ratio}
\label{sec:ratio_analysis}

Figure~\ref{fig:ratio_effect} (described textually below) illustrates the
degradation profile as training data shifts from 1:1 to 5:1.

\textbf{Two-tower degradation mechanism.}
The two-tower model's prediction is $\sigma(\langle \mathbf{z}^{\text{TF}},
\mathbf{z}^{\text{Gene}} \rangle / \tau)$.
At training time with 5:1 negatives, the model optimizes BCE loss on a dataset
where the class prior is $p(\text{pos}) = 1/6$.
The optimal decision boundary for $\sigma(\cdot) = 0.5$ corresponds to equal
positive and negative posterior probability.
Under a 5:1 class ratio, the model's calibrated threshold should be
$\sigma^{-1}(1/6) \approx -1.61$ logits rather than 0.
Since we evaluate at threshold $= 0.5$ (logit $= 0$) on a 1:1 test set,
all pairs with true posterior between $1/6$ and $1/2$ will be misclassified.
The two-tower's cosine-similarity scoring function provides limited control
over this calibration: the score distribution is coupled to the angular
structure of the embedding space, making it difficult to independently shift
the operating threshold without changing the embedding geometry.

\textbf{Cross-encoder robustness.}
The cross-encoder outputs a raw logit whose scale is unconstrained.
When trained with 5:1 negatives, the network can learn to shift its logit
distribution toward more negative values, while the absolute operating point
(threshold $= 0$ after sigmoid) can be compensated by the bias term in the
final layer.
This decoupling of the learned representation from the scoring calibration
makes the cross-encoder substantially more robust to negative-ratio shift.

\begin{table}[!htbp]
\centering
\caption{Accuracy degradation from 1:1 to 5:1 training. $\Delta$Acc: absolute
accuracy drop (test set always 1:1).}
\label{tab:ratio_degradation}
\small
\begin{tabular}{@{}lcccc@{}}
\toprule
\textbf{Model} & \textbf{1:1 Acc.} & \textbf{5:1 Acc.} & \textbf{$\Delta$Acc} & \textbf{$\Delta$AUROC} \\
\midrule
Two-Tower     & 80.90\% & 74.03\% & $-6.87$~pp & $-0.067$ \\
Cross-Encoder & 83.03\% & 81.47\% & $-1.56$~pp & $+0.011$ \\
\bottomrule
\end{tabular}
\end{table}

Strikingly, the cross-encoder's AUROC \emph{increases} by $+0.011$ when
moving to 5:1 training (Table~\ref{tab:ratio_degradation}).
AUROC measures ranked discrimination independent of threshold, so a model
that trains on harder negatives may learn sharper score separation even if its
accuracy at a fixed threshold is marginally lower.
The cross-encoder at 5:1 achieves AUROC = 0.915---the highest of any model
variant in this study---suggesting that realistic negative sampling provides
richer training signal for the monolithic scoring function.

\subsection{Ensemble performance}
\label{sec:ensemble_results}

Ensemble averaging (5 models) provides consistent gains:
Two-tower 1:1: $80.90\% \to 83.47\%$ ($+2.57$~pp);
Cross-encoder 1:1: $83.03\% \to 84.02\%$ ($+0.99$~pp).
The smaller ensemble gain for the cross-encoder reflects lower individual model
variance and higher correlation among ensemble members.
At 5:1, the cross-encoder ensemble (82.33\%) outperforms the two-tower ensemble
(75.92\%) by 6.41~pp.

\subsection{Cross-seed stability}
\label{sec:stability}

The cross-encoder exhibits lower cross-seed variance at 1:1 (SD = 0.48~pp
vs.\ 0.59~pp for two-tower) and 5:1 (0.74~pp vs.\ 1.54~pp).
At 5:1, the two-tower's worst-performing seed (789) achieves only 71.11\%
accuracy, a 4.25~pp drop below the mean---nearly double the cross-encoder's
maximum deviation (1.01~pp).
This stability advantage likely reflects the cross-encoder's unconstrained
logit range, which provides a wider gradient signal corridor during early
training regardless of initialization.

% ============================================================
\section{Discussion}
\label{sec:discussion}

\subsection{Why does the cross-encoder outperform the two-tower?}

The cross-encoder's advantage on this dataset likely stems from two structural
properties:

\textbf{Explicit interaction features.}
The element-wise product $\mathbf{e}^{\text{TF}} \odot \mathbf{e}^{\text{Gene}}$
directly encodes multiplicative embedding interactions.
In the two-tower model, the dot product $\langle \mathbf{z}^{\text{TF}},
\mathbf{z}^{\text{Gene}} \rangle$ implicitly computes a \emph{sum} of these
products after nonlinear transformation, but the intermediate layers cannot
access raw multiplicative features.
For GRN edge prediction, where binding affinity is plausibly related to
multiplicative compatibility of TF binding domain and gene promoter
features~\cite{yang2015embedding}, explicit interaction features may provide a
direct computational path that the two-tower must reconstruct expensively.

\textbf{Unconstrained logit vs.\ fixed sigmoid range.}
The two-tower produces a sigmoid output constrained to $(0, 1)$, with the score
distribution's shape determined by the geometry of the learned embedding space
and the fixed temperature $\tau = 0.05$.
The cross-encoder outputs a raw logit in $(-\infty, +\infty)$, giving the
optimizer full control over score calibration.
This matters especially under class-ratio shift (\S\ref{sec:ratio_analysis}):
the cross-encoder can adapt its output distribution by adjusting the final-layer
bias, while the two-tower must restructure its embedding geometry.

\textbf{Why the AUROC gap is larger than the accuracy gap.}
AUROC (0.904 vs.\ 0.810, $\Delta = 0.094$) is proportionally much larger than
the accuracy gap (83.03\% vs.\ 80.90\%, $\Delta = 2.13$~pp).
AUROC measures the quality of the full score ranking, not just the binary
decision at threshold 0.5.
The two-tower's temperature-scaled cosine similarity produces a compact score
distribution: because embeddings are normalized through FC layers but the
dot product is unscaled before temperature division, the effective score range
is narrow.
The cross-encoder's logit-scale output produces a wider, better-separated score
distribution, yielding stronger ranking quality even at modest accuracy gains.

\subsection{Practical implications of gradient failure analysis}

The three failure modes identified in \S\ref{sec:grad_analysis} are each
individually capable of collapsing training and are easily missed:

\begin{itemize}
    \item \textbf{Double-sigmoid}: arises whenever a utility function
    re-applies sigmoid to an already-sigmoid output, a common mistake when
    re-using loss utilities with architectures that differ in output activation.
    \item \textbf{Gradient replacement}: a subtle difference between $\mathbf{g}
    \mathrel{+}= \nabla$ (correct for epoch-level accumulation) and $\mathbf{g}
    \mathrel{=} \nabla$ (the implementation) that silently reduces the effective
    batch size to a single mini-batch.
    \item \textbf{Clip-temperature interaction}: the effective gradient norm
    scales as $1/\tau$, so a clip threshold tuned for $\tau = 1.0$ may be
    100$\times$ too small for $\tau = 0.01$, annihilating all weight updates.
\end{itemize}

These failures are relevant to any from-scratch neural network implementation
in systems languages (Rust, C++) that lack automatic gradient checking.
We recommend: (1) verify that loss utility functions do not apply activations
redundantly; (2) prefer per-batch optimizer steps over epoch-level accumulation
when backward functions replace (rather than accumulate) stored gradients; and
(3) scale clip thresholds relative to $1/\tau$ when using temperature-scaled
scoring functions.

\subsection{Architectural trade-offs and when to prefer each design}

\textbf{Prefer the two-tower model} when:
\begin{itemize}
    \item Large-scale retrieval is needed (e.g., ranking all $M \times N$ pairs):
    tower outputs $\mathbf{z}^{\text{TF}}$ and $\mathbf{z}^{\text{Gene}}$ can
    be pre-computed and dot-product search applied via FAISS.
    \item Interpretable latent geometry is desired: TF and gene positions in
    the shared $\mathbb{R}^{512}$ space encode entity-level properties
    independent of query.
    \item Training is class-balanced (1:1): the accuracy and F1 gap to the
    cross-encoder is only $\sim$2~pp.
\end{itemize}

\textbf{Prefer the cross-encoder} when:
\begin{itemize}
    \item Training data is imbalanced (5:1 or higher): the cross-encoder
    maintains strong performance across negative ratios, while the two-tower
    degrades substantially.
    \item AUROC is the primary metric: the cross-encoder achieves 0.904--0.915
    vs.\ 0.743--0.810 for the two-tower.
    \item Downstream evaluation requires score calibration (e.g., top-$k$
    ranked lists for experimental validation).
\end{itemize}

\subsection{Limitations}
\label{sec:limitations}

Several limitations bound the scope of these conclusions.
First, the current evaluation uses 0\% expression coverage (all entity expression
vectors are zero-padded), meaning both models operate exclusively on learnable
embeddings.
Incorporating non-zero expression features may shift the relative advantage.
Second, evaluation is limited to a single tissue (human brain) and two
databases; generalization to other tissues, organisms, or regulatory contexts
requires further validation.
Third, the test set evaluation uses 1:1 class balance regardless of training
ratio; a held-out test set matching the training ratio distribution would provide
a complementary perspective.
Fourth, we do not evaluate cold-start generalization (new TFs or genes unseen
in training), which is a harder and more biologically relevant setting.
Fifth, comparisons with GPU-based transformer methods remain approximate due to
dataset and protocol differences; a shared benchmark such as
BEELINE~\cite{pratapa2020benchmarking} is needed for rigorous comparison.

% ============================================================
\section{Conclusion}
\label{sec:conclusion}

We presented a controlled comparison of modular (two-tower) and monolithic
(cross-encoder) neural architectures for gene regulatory network edge prediction.
A prerequisite diagnostic investigation identified three compounding gradient
failures in the original two-tower training procedure---double-sigmoid in the
backward pass, per-batch gradient replacement in epoch-level updates, and
clip-threshold incompatibility with temperature scaling---each sufficient to
collapse training to near-chance accuracy.
After resolving these failures, the corrected two-tower achieves 80.90\%
accuracy (AUROC 0.810) under balanced training, and the cross-encoder achieves
83.03\% accuracy (AUROC 0.904) under identical conditions.

The cross-encoder's advantage is especially pronounced in AUROC ($+0.094$ at
1:1) and in robustness to negative-ratio shift: cross-encoder accuracy drops
only 1.56~pp at 5:1 training versus 6.87~pp for the two-tower.
We attribute this to the cross-encoder's explicit interaction features
($\mathbf{e}^{\text{TF}} \odot \mathbf{e}^{\text{Gene}}$) and unconstrained
logit output, which together provide richer interaction representation and
better calibration control.

For practitioners, we recommend: (1) implement gradient sanity checks in
from-scratch implementations; (2) prefer per-batch optimizer steps over
epoch-level updates when backward functions do not accumulate gradients;
(3) use the cross-encoder design when training data is imbalanced or AUROC
is the primary metric; (4) reserve the two-tower for large-scale retrieval
settings where pre-computed tower representations provide computational efficiency.

\subsection*{Acknowledgments}
Computations were performed on consumer-grade Apple M-series hardware.
We thank the developers of DoRothEA and TRRUST for curated regulatory data.

\subsection*{Author contributions}
E.L.\ conceived the study, designed and implemented both architectures,
performed the gradient analysis and all experiments, and wrote the manuscript.

\subsection*{Data availability}
Source code, result JSONs, trained model weights, and evaluation scripts are
available at \url{https://github.com/Evintkoo/module-regularized-grn}.

% ============================================================

\bibliographystyle{unsrtnat}
\bibliography{references}

% ============================================================
\clearpage
\onecolumn
\appendix
\section*{Supplementary Material}
\addcontentsline{toc}{section}{Supplementary Material}

\subsection*{S1. Full hyperparameter specification}

\begin{table}[H]
\centering
\caption*{\textbf{Table S1.} Hyperparameter configuration for both models.}
\small
\begin{tabular}{@{}lllcc@{}}
\toprule
\textbf{Category} & \textbf{Parameter} & & \textbf{Two-Tower} & \textbf{Cross-Encoder} \\
\midrule
\multirow{3}{*}{Architecture}
  & Embedding dim.\ ($d_e$)   & & 512 & 512 \\
  & Hidden dim.\ ($d_h$)      & & 512 & 512 \\
  & Output/logit dim.         & & 512 (sigmoid) & 1 (logit) \\
\midrule
\multirow{4}{*}{Optimizer (Adam)}
  & Learning rate ($\eta$)    & & $10^{-3}$ & $10^{-3}$ \\
  & $\beta_1$                 & & 0.9 & 0.9 \\
  & $\beta_2$                 & & 0.999 & 0.999 \\
  & Gradient clip             & & 5.0 & 5.0 \\
\midrule
\multirow{2}{*}{Training}
  & Batch size                & & 256 & 256 \\
  & Epochs                    & & 60 & 60 \\
\midrule
Scoring
  & Temperature ($\tau$)      & & 0.05 & N/A \\
\midrule
\multirow{2}{*}{Initialization}
  & Embeddings $\sigma$       & & 0.01 & 0.10 \\
  & FC layers                 & & $\sigma=0.01$ & He init \\
\bottomrule
\end{tabular}
\end{table}

\subsection*{S2. Per-seed results: 5:1 negative sampling}

\begin{table}[H]
\centering
\caption*{\textbf{Table S2.} Per-seed results under 5:1 negative sampling (test set 1:1 balanced).}
\small
\begin{tabular}{@{}cccccccc@{}}
\toprule
\multirow{2}{*}{\textbf{Seed}} &
\multicolumn{3}{c}{\textbf{Two-Tower (5:1)}} & &
\multicolumn{3}{c}{\textbf{Cross-Encoder (5:1)}} \\
\cmidrule{2-4} \cmidrule{6-8}
 & Acc.\ (\%) & AUROC & F1 & & Acc.\ (\%) & AUROC & F1 \\
\midrule
42   & 73.93 & 0.741 & 0.659 & & 81.64 & 0.923 & 0.785 \\
123  & 75.17 & 0.757 & 0.680 & & 80.70 & 0.916 & 0.769 \\
456  & 75.36 & 0.753 & 0.684 & & 81.99 & 0.924 & 0.790 \\
789  & 71.11 & 0.714 & 0.609 & & 82.47 & 0.913 & 0.797 \\
1337 & 74.57 & 0.752 & 0.668 & & 80.55 & 0.899 & 0.771 \\
\midrule
Mean $\pm$ SD
     & 74.03 $\pm$ 1.54 & 0.743 & 0.660
     & & 81.47 $\pm$ 0.74 & 0.915 & 0.783 \\
\bottomrule
\end{tabular}
\end{table}

\subsection*{S3. Comparison with Version 1 results}

Version 1 of this work reported a Two-Tower single model with 80.14\% accuracy
and a 5-model ensemble with 83.06\% accuracy.
These figures correspond to a different training implementation using epoch-level
SGD and the \texttt{bce\_loss\_backward} utility; they may reflect the
interaction between Failure~1 (double-sigmoid) and the epoch-level gradient
accumulation, which in combination produced partially functional training.
Version 2 uses per-batch Adam with the stable gradient estimator (Eq.~\ref{eq:stable_grad}),
achieving 80.90\% $\pm$ 0.59\% (mean over 5 seeds) and 83.47\% ensemble,
which are consistent with V1 figures to within the seed variance.
The V2 training procedure is strictly more principled: gradient signals are
numerically stable, all parameters update every batch, and the optimizer state
is fully defined.

\end{document}
